\documentclass[11pt]{article}
\usepackage{amsmath,amssymb}
\usepackage[margin=1in]{geometry}

\title{Blueprint for the Ternary Quartic Rank-4 Proof}
\date{2026-04-09}

\begin{document}
\maketitle

\section*{Status}
This blueprint is not complete yet. The abstract SOCP component has been
formalized, but the ternary-quartic algebraic reduction is still missing.

\section*{Formalized Step}
Let $r(u,p) = \sigma(u) - p$. The following statement is already formalized in
\texttt{TernaryQuarticProof/Certificate.lean}.

\medskip
\noindent
\textbf{Lemma.} Assume:
\begin{enumerate}
\item $u$ is an admissible rank-4 point;
\item $u$ is a SOCP for $p$ with respect to a positive-definite bilinear form
$B$;
\item $r(u,p)$ is orthogonal to $\operatorname{im}(A_{u_{\mathrm{img}}})$ for
some admissible factor $u_{\mathrm{img}}$;
\item there is a decomposition
\[
  p = q_{\mathrm{img}} + \sum_{w \in W} \sigma(w),
\]
where $q_{\mathrm{img}} \in \operatorname{im}(A_{u_{\mathrm{img}}})$ and every
$w \in W$ is an admissible kernel direction for $A_u$.
\end{enumerate}
Then $r(u,p) = 0$.

\medskip
\noindent
\textbf{Proof.}
The image term pairs to zero by the first-order condition. Each kernel-square
term pairs nonnegatively with the residual by the second-order condition,
because $A_u(w)=0$ removes the positive $2\langle A_u(w), A_u(w)\rangle$ part
of the Hessian expression. Summing these inequalities gives
$\langle p, r(u,p)\rangle \ge 0$. On the other hand, the first-order condition
at $v=u$ gives $\langle \sigma(u), r(u,p)\rangle = 0$, hence
\[
  \langle r(u,p), r(u,p)\rangle
  = \langle \sigma(u)-p, r(u,p)\rangle
  = - \langle p, r(u,p)\rangle \le 0.
\]
Positive-definiteness implies $\langle r(u,p), r(u,p)\rangle \ge 0$, so the
objective is zero and therefore $r(u,p)=0$.

\section*{Additional Formalized Reductions}
Two special cases are now formalized.

\medskip
\noindent
\textbf{Image case.} If $p \in \operatorname{im}(A_u)$ with an admissible
preimage, then $r(u,p)=0$ directly from the first-order condition and positive
definiteness.

\medskip
\noindent
\textbf{Constant-relation case.} If there exists a nonzero constant vector
$c \in \mathbb{R}^4$ with $\sum_i c_i u_i = 0$, then $r(u,p)=0$ for every SOS
quartic $p$.

\medskip
\noindent
\textbf{Proof sketch.}
For each quadratic summand $q^2$ in an SOS decomposition of $p$, define the
admissible direction $w_i = c_i q$. Then
\[
  A_u(w) = q \sum_i c_i u_i = 0,
  \qquad
  \sigma(w) = \Bigl(\sum_i c_i^2\Bigr) q^2.
\]
Since $\sum_i c_i^2 > 0$, the SOCP inequality in the direction $w$ implies
$\langle q^2, r(u,p)\rangle \ge 0$. Summing over the SOS decomposition yields
$\langle p, r(u,p)\rangle \ge 0$, and the abstract SOCP argument finishes as
before.

\section*{Normal-Form Target}
The most promising unresolved case is when, after affine changes of variables
and invertible recombination of the four coordinates, the factor takes the form
\[
  u = (1, x_1, x_2, q).
\]
For this normal form there is an explicit admissible kernel family:
if a homogeneous quadratic $h$ is written as
\[
  h = x_1 a + x_2 b
\]
with $a,b$ linear, then
\[
  w = (-h, a, b, 0)
\]
satisfies $A_u(w)=0$ and
\[
  \sigma(w) = h^2 + a^2 + b^2.
\]
The correction term $a^2+b^2$ has degree at most $2$, hence lies in the image
through the first coordinate $1$. Therefore every square of a homogeneous
quadratic belongs to
\[
  \operatorname{im}(A_u) + \operatorname{cone}(\sigma(\ker A_u)).
\]
The remaining gap is to turn this observation into a complete classification of
all independent rank-$4$ factors of image rank $13$ or $14$.

\section*{Pending Step}
To finish the theorem, it remains to prove the multivariate admissible
image-plus-cone decomposition for every SOS quartic. This is currently the main
mathematical gap. The likely form is a finite-dimensional quadratic-space
certificate rather than a direct copy of the univariate gcd/peeling argument.

\section*{Planned File Map}
\begin{itemize}
\item \texttt{TernaryQuarticProof/Socp.lean}: abstract SOCP consequences.
\item \texttt{TernaryQuarticProof/Certificate.lean}: admissible transformed
image-plus-cone certificate.
\item one future algebra module for the ternary-quartic decomposition theorem;
\item \texttt{TernaryQuarticProof.lean}: final theorem assembly.
\end{itemize}

\end{document}
